\documentclass[11pt]{article}
\usepackage[margin=1in]{geometry}
\usepackage{booktabs}
\usepackage{graphicx}
\usepackage{float} % Keep each figure near its discussion
\usepackage{placeins} % Prevent floats from crossing section boundaries
\usepackage{url}
\usepackage[hidelinks]{hyperref} % keep last

\begin{document}

\input{front_page}

\tableofcontents
\clearpage



\section{Problem statement}

Graph-based approximate nearest-neighbor (ANN) search balances recall,
latency, and index size. SHG (VLDB 2025)~\cite{shg} reports better recall--latency
trade-offs than HNSW~\cite{hnsw} (Figure~\ref{fig:paper-curve}). We tested this claim using the released SHG code~\cite{shgcode} and a native HNSW baseline (hnswlib~\cite{hnswlib}).

\begin{figure}[H]
\centering
\includegraphics[width=\linewidth,height=0.34\textheight,keepaspectratio]{figures/shg_paper_recall_latency.png}
\caption{Recall--latency results reported by the SHG authors~\cite{shg} (their Figure 5).}
\label{fig:paper-curve}
\end{figure}

We evaluated Enron~\cite{survey}, GIST1M, and SIFT1M~\cite{pq,texmex} ($M=48$,
$ef_{construction}=80$, $k=20$, $ef_{search}=20$--512).
HNSW was faster at comparable recall on Enron and SIFT1M. On GIST1M,
SHG did not reach 0.90 recall, while HNSW reached 0.934.
SHG also used more index memory (Table~\ref{tab:comparison}).

\begin{table}[H]
\centering
\small
\begin{tabular}{lcccc}
\toprule
Dataset & \multicolumn{2}{c}{Latency at $\geq$0.90 recall ($\mu$s)}
        & Max. recall@$k=20$ & Index size (MiB) \\
        & HNSW & SHG & HNSW / SHG & HNSW / SHG \\
\midrule
Enron  & 416  & 808         & 0.9925 / 0.9910 & 533 / 698 \\
GIST1M & 3308 & Not reached & 0.9632 / 0.6300 & 4048 / 5249 \\
SIFT1M & 199  & 400         & 0.9988 / 0.9982 & 874 / 1026 \\
\bottomrule
\end{tabular}
\caption{Mean full-query latency at the first measured sweep point reaching
0.90 recall, maximum recall across the sweep, and saved index size.}
\label{tab:comparison}
\end{table}

Our work instead targets memory stalls without changing search decisions.

\begin{figure}[H]
\centering
\includegraphics[width=\linewidth,height=0.34\textheight,keepaspectratio]{figures/recall_latency_vanilla_all_datasets.png}
\caption{Measured recall--latency curves for the three datasets for M=48 and efConstruction=80.}
\label{fig:our-curve}
\end{figure}

\section{Reproduction limitations}

This is not an exact reproduction: we use SIFT1M rather than SIFT100M~\cite{texmex},
and a separate HNSW implementation~\cite{hnswlib}. SHG also uses compressed distances
at build time. Our curves sweep $ef_{search}$ over all queries, whereas
the paper uses recall-ordered subsets; its plotting script is unavailable.

The released example lacks a working vanilla HNSW baseline and uses
$k$, not $ef_{search}$, as SHG's base-layer breadth. We use
$\max(ef_{search},k)$. SHG skipped no levels because its shortcut tables
were below the 100-entry activation threshold. These differences limit
comparison with the published results.

\section{Experimental setup}

We run our experiments with one thread on an Ubuntu 22.04 machine equipped with a 24-core Intel Xeon Gold 5220R CPU at 2.20 GHz and 125 GiB of RAM. Base-layer profiling used $ef_{search}=64$.

\section{Observations}
\label{sec:observation}
\subsection{Dependent memory accesses}

Each expanded node requires a neighbor-list read, visited-flag checks,
and vector loads for unseen neighbors. These dependent accesses can stall
the core.

\subsection{Where query time is spent}

Figures~\ref{fig:hnsw-time-breakdown} and~\ref{fig:shg-time-breakdown}
break query time into distance computation, memory-related idle time,
and search overhead.

\begin{figure}[H]
\centering
\includegraphics[width=\linewidth,height=0.34\textheight,keepaspectratio]{figures/21_full_breakdown_hnsw_ef64.pdf}
\caption{Breakdown of average base-layer query time for HNSW at
$ef_{search}=64$.}
\label{fig:hnsw-time-breakdown}
\end{figure}

\begin{figure}[H]
\centering
\includegraphics[width=\linewidth,height=0.34\textheight,keepaspectratio]{figures/21b_full_breakdown_shg_ef64.pdf}
\caption{Breakdown of average base-layer query time for SHG at
$ef_{search}=64$.}
\label{fig:shg-time-breakdown}
\end{figure}

\begin{table}[H]
\centering
\small
\begin{tabular}{lcccc}
\toprule
 & \multicolumn{2}{c}{HNSW} & \multicolumn{2}{c}{SHG} \\
Dataset & Avg. time ($\mu$s) & Idle time & Avg. time ($\mu$s) & Idle time \\
\midrule
Enron  & $\sim$346 & 62\,$\mu$s (18\%) & $\sim$451 & 75\,$\mu$s (17\%) \\
GIST1M & $\sim$495 & 68\,$\mu$s (14\%) & $\sim$299 & 70\,$\mu$s (24\%) \\
SIFT1M & $\sim$94 & 35\,$\mu$s (37\%) & $\sim$181 & 43\,$\mu$s (24\%) \\
\bottomrule
\end{tabular}
\caption{Base-layer query time and memory-related idle share at
$ef_{search}=64$; totals are inferred from rounded values.}
\label{tab:query-time-breakdown}
\end{table}

On SIFT1M, HNSW spends about 37\% of its 0.103\,ms base-layer time idle;
SHG spends 24\% of 0.196\,ms idle. Across datasets, idle shares range
from 13--37\% for HNSW and 16--24\% for SHG. This motivates early,
selective data requests rather than indiscriminate prefetching.

\subsection{Why existing prefetching falls short}

Native HNSW~\cite{hnswlib} prefetches the next neighbor's flag and first vector cache
line, potentially too late to hide latency. It may also fetch already-visited
vectors. The released SHG base-layer loop~\cite{shgcode} has no software prefetches.

Whole-vector prefetching slowed Enron and GIST1M by 10--20\% versus
no prefetching. Unfiltered look-ahead helped SIFT1M but slowed HNSW by
1.5\% on Enron and 5.8\% on GIST1M. Extra traffic can outweigh the benefit.

\section{Optimization direction}

Memory stalls account for substantial query time, but indiscriminate
prefetching can be slower. Our goal is to \emph{prefetch useful data early
without excess memory traffic}, while preserving search behavior.

\section{Solution: Pf\_Filter}
\label{sec:pf-filter}

\textsc{Pf\_Filter} looks ahead in the base-layer neighbor loop and
uses visited flags to avoid prefetching unnecessary vectors. It changes only when data is loaded, not candidate selection or distance computation.

\subsection{Two-stage filtered look-ahead}

The visited flag determines whether a vector is needed. With look-ahead
$D=2$ (Figure~\ref{fig:pf-filter-design}), processing neighbor $j$:
\begin{enumerate}
  \item Prefetches the visited flag of neighbor $j+2D$ ($j+4$).
  \item Checks the flag of neighbor $j+D$ ($j+2$), which was prefetched two neighbours earlier and is normally in cache, and prefetches the first three cache lines of its vector only if that neighbour has not been visited.
\end{enumerate}

This staggers dependent loads and avoids speculative requests for visited
vectors.

\subsection{Preparing the next accesses}

At the start of each neighbour list, Pf\_Filter prefetches the first $2D$ visited flags and the first $D$ vectors, so the look-ahead is full from the first neighbour. It also prefetches the neighbour list of the next candidate right after a pop, and that of the queue's top node after each push. SHG additionally prefetches the target of a shortcut jump; this part is not used in HNSW.

\subsection{Correctness and overhead}

Prefetches do not change search state or index layout. Tested
configurations had zero result mismatches against vanilla HNSW; overhead is
limited to extra flag checks and prefetch instructions.


\begin{figure}[H]
\centering
\includegraphics[width=\linewidth,height=0.34\textheight,keepaspectratio]{figures/pf_filter_steps.pdf}
\caption{Two-stage prefetching with $D=2$: while processing neighbor $j$,
\textsc{Pf\_Filter} requests the flag for $j+4$ and conditionally requests
vector cache lines for $j+2$.}
\label{fig:pf-filter-design}
\end{figure}




\section{Results}
\label{sec:results}

Figure~\ref{fig:pf-filter-results} shows that \textsc{Pf\_Filter} shifts
recall--latency curves left: lower latency at unchanged recall for both
algorithms. Each point averages five full-query-set runs ($M=48$,
$ef_{construction}=80$, $k=20$).

\begin{figure}[H]
\centering
\includegraphics[width=\linewidth,height=0.34\textheight,keepaspectratio]{figures/recall_latency_all_datasets.png}
\caption{Recall--latency plot for Enron, GIST1M,
and SIFT1M for M=48 and efConstruction=80.}
\label{fig:pf-filter-results}
\end{figure}

At $ef_{search}=64$, HNSW latency falls by 20.0\% (Enron), 24.3\%
(GIST1M), and 25.6\% (SIFT1M); SHG falls by 17.7\%, 26.2\%, and
18.5\%, respectively. SHG's recall deficit on GIST1M remains.


\subsection{Throughput}

As $ef_{search}$ rises, latency increases throughput decreases (Figure~\ref{fig:pf-filter-qps}).
\textsc{Pf\_Filter} consistently improves both metrics over each
algorithm's baseline.


\begin{figure}[H]
\centering
\includegraphics[width=\linewidth,height=0.34\textheight,keepaspectratio]{figures/qps_all_datasets.png}
\caption{Search throughput across $ef_{search}$ values. Higher QPS indicates
more queries completed per second.}
\label{fig:pf-filter-qps}
\end{figure}

Across ten $ef_{search}$ values, filtered HNSW averages
$1.33\times$ speedup on SIFT1M, $1.27\times$ on Enron, and
$1.33\times$ on GIST1M (Figure~\ref{fig:pf-filter-speedup}).
Ratios are normalized to vanilla HNSW at equal $ef_{search}$;
recall may differ.

\begin{figure}[H]
\centering
\includegraphics[width=\linewidth,height=0.34\textheight,keepaspectratio]{figures/mean_latency_speedup_all_datasets.png}
\caption{Mean same-$ef_{search}$ speedup over vanilla HNSW;
recall may differ.}
\label{fig:pf-filter-speedup}
\end{figure}

\subsection{Effect of Pf\_Filter on Base-Layer Time}
\label{sec:pf-filter-breakdown}

Figures~\ref{fig:pf-filter-breakdown-hnsw} and~\ref{fig:pf-filter-breakdown-shg}
show where the base-layer time goes at $ef_{search}=64$, before and after adding
\textsc{Pf\_Filter}. The base layer takes 93--99.8\% of the query time as discussed in section~\ref{sec:observation}.

\textsc{Pf\_Filter} cuts the idle time on every dataset, by 10--15\% for HNSW and
16--21\% for SHG. Distance computation does not change. For HNSW the total saving is
small (0.5--2.6\%), because the search overhead stays about the same. For SHG the
overhead also drops, so the base layer is 8--15\% faster. The search results do not
change.

\begin{figure}[H]
\centering
\includegraphics[width=\linewidth,height=0.34\textheight,keepaspectratio]{figures/19a_problem_vs_pf_filter_hnsw.pdf}
\caption{Base-layer time per query for vanilla HNSW (left bar) and HNSW +
\textsc{Pf\_Filter} (right bar) at $ef_{search}=64$. Above each pair: change in total
time and in idle time.}
\label{fig:pf-filter-breakdown-hnsw}
\end{figure}

\begin{figure}[H]
\centering
\includegraphics[width=\linewidth,height=0.34\textheight,keepaspectratio]{figures/19b_problem_vs_pf_filter_shg.pdf}
\caption{Base-layer time per query for SHG as published (left bar) and SHG +
\textsc{Pf\_Filter} (right bar) at $ef_{search}=64$. Above each pair: change in total
time and in idle time.}
\label{fig:pf-filter-breakdown-shg}
\end{figure}


\clearpage
\phantomsection
\addcontentsline{toc}{section}{References}
\begin{thebibliography}{9}

\bibitem{shg}
Z.~Gong, Y.~Zeng, and L.~Chen.
\newblock Accelerating approximate nearest neighbor search in hierarchical graphs:
  Efficient level navigation with shortcuts.
\newblock \emph{PVLDB}, 18(10):3518--3530, 2025.
\newblock doi:10.14778/3748191.3748212.

\bibitem{shgcode}
Z.~Gong.
\newblock SHG-Index: source code.
\newblock \url{https://github.com/gzyhkust/SHG-Index}.

\bibitem{hnsw}
Y.~A. Malkov and D.~A. Yashunin.
\newblock Efficient and robust approximate nearest neighbor search using
  hierarchical navigable small world graphs.
\newblock \emph{IEEE Trans. Pattern Anal. Mach. Intell.}, 42(4):824--836, 2020.

\bibitem{hnswlib}
hnswlib: header-only C++/Python library for fast approximate nearest neighbors.
\newblock \url{https://github.com/nmslib/hnswlib}.

\bibitem{pq}
H.~J\'egou, M.~Douze, and C.~Schmid.
\newblock Product quantization for nearest neighbor search.
\newblock \emph{IEEE Trans. Pattern Anal. Mach. Intell.}, 33(1):117--128, 2011.

\bibitem{texmex}
H.~J\'egou, M.~Douze, and C.~Schmid.
\newblock Datasets for approximate nearest neighbor search (TEXMEX corpus).
\newblock \url{http://corpus-texmex.irisa.fr/}.

\bibitem{survey}
M.~Wang, X.~Xu, Q.~Yue, and Y.~Wang.
\newblock A comprehensive survey and experimental comparison of graph-based
  approximate nearest neighbor search.
\newblock \emph{PVLDB}, 14(11):1964--1978, 2021.

\end{thebibliography}

\end{document}